
Cuando Morse desarrolló su teoría, a lo largo de los años veinte (\cite{BottMarstonMorse}), la escena matemática era muy distinta, no fue hasta $1949$ que Whitehead formalizó siquiera el concepto de CW-complejo tal y como lo conocemos (\cite{KletteCellComplexesHistory}). En particular, el teorema \ref{cwvariedad} aún no había sido demostrado, por lo que la relación entre la topología de una variedad diferenciable y los puntos críticos de una función de Morse fue descrita mediante una colección de desigualdades. 

En este capítulo presentamos esta formulación original siguiendo de nuevo el desarrollo encontrado en \cite{milnor1963}. 

\section{Subaditividad}

Empezamos introduciendo algunos resultados de funciones subaditivas.

A lo largo de este capítulo nos referiremos a pares de espacios $(X,Y)$. Conviene aclarar que no se trata de pares de espacios arbitrarios, sino de pares tales que $Y$ es un subespacio de $X$ para el que tiene sentido construir los grupos de homología relativa (ver \ref{sec:hom_relat}). Recordemos también que en el caso en el que $Y$ es vacío o contractible recuperamos la homología de $X$.

\begin{definition}
    Sea $S$ una función que toma pares $(X,Y)$ de espacios topológicos y les asigna un entero $S(X,Y)\in \mathbb{Z}$. Decimos que $S$ es \textit{subaditiva} si siempre que $X\supset Y\supset Z$ se tiene que
    \begin{equation*}
        S(X,Z) \le S(X,Y) + S(Y,Z).
    \end{equation*}
    Si en las misma hipótesis se verifica siempre la igualdad diremos que $S$ es \textit{aditiva}.
\end{definition}

\begin{example}
    \hfill
    \begin{enumerate}[label=(\roman*)]
        \item La función $b_k(X,Y)$ que asigna a cada par $(X,Y)$ su $k$-ésimo número de Betti es subaditiva, siempre que sean finitos. En efecto, sean $X$, $Y$ y $Z$ tales que $X\supset Y\supset Z$, se tiene la siguiente secuencia exacta:
        \begin{align*}
            \cdots \longrightarrow H_k(Y,Z) \overset{f}{\longrightarrow} H_k(X,Z) \overset{g}{\longrightarrow} H_k(X,Y) \overset{h}{\longrightarrow} \cdots
        \end{align*}
        Esta surge de generalizar la secuencia mencionada en \ref{exact_sec_hom_rel} a una terna (\cite[p.~118]{hatcher2002}). De ella se deduce la subaditividad, veamos cómo. Dado un homomorfismo entre grupos abelianos de rango finito $\varphi: G\to H$, por el primer teorema de isomorfía $G/\text{ker}\varphi$ es isomorfo a $\text{im} \varphi$, luego se tiene
        \begin{equation*}
            \text{rank } G = \text{rank}(\text{ker } \varphi) + \text{rank}(\text{im } \varphi).
        \end{equation*}
        En nuestro caso, tenemos las siguientes igualdades
        \begin{equation*}
            \text{im } f = \text{ker } g ~~\text{ y }~~ \text{im } g = \text{ker } h.
        \end{equation*}
        Además,
        \begin{align*}
            &b_k(Y,Z) = \text{rank}(\text{ker } f) + \text{rank}(\text{im } f) = \text{rank}(\text{ker } f) + \text{rank}(\text{ker } g),\\
            &b_k(X,Z) = \text{rank}(\text{ker } g) + \text{rank}(\text{im } g),\\
            &b_k(X,Y) = \text{rank}(\text{ker } h) + \text{rank}(\text{im } h) = \text{rank}(\text{im } g) + \text{rank}(\text{im } h). 
        \end{align*}
        Por lo que
        \begin{equation*}
            b_k(X,Z) - b_k(Y,Z) = \text{rank}(\text{im } g) - \text{rank}(\text{ker } f) \le b_k(X,Y).
        \end{equation*}
        De donde se deduce la subaditividad.
        \item Aunque la demostración es más involucrada y no la añadiremos, se tiene que la funcíon que asigna a cada par su característica de Euler-Poincaré $\chi (X,Y)$ es aditiva.
    \end{enumerate}
\end{example}

\bigskip

\begin{lemma}
    Sea $S$ una función subaditiva y sean $X_0\subset \dots \subset X_n$. Entonces,
    \begin{equation*}
        S(X_n, X_0) \le \sum^n_{i=1} S(X_i,X_{i-1}).
    \end{equation*}
    Si $S$ es aditiva se verifica la igualdad.
\end{lemma}
\begin{proof}
    Procedemos por inducción sobre $n$. Para $n=1$ se verifica la igualdad trivialmente y para $n=2$ se tiene la definición de subaditividad.


    
\end{proof}

\bigskip

\begin{lemma}
    La función $S_{\lambda}$ es subaditiva, donde 
    \begin{equation*}
        S_{\lambda}(X,Y) = \sum_{i=0}^{\lambda} (-1)^i b_{\lambda-i}(X,Y)
    \end{equation*}
\end{lemma}
\begin{proof}

\end{proof}

\bigskip

%%%%%%%%%%%%%%%%%%%%%%%%%%%%%%%%%%%%
%%                                %%
%%       LAS DESIGUALDADES        %%
%%                                %%
%%%%%%%%%%%%%%%%%%%%%%%%%%%%%%%%%%%%
\section{Las desigualdades}

\begin{theorem}[Desigualdades de Morse, versión débil]
    Sean $f$ una función suave {¿\color{blue} de Morse?} sobre una variedad compacta M y $C_{\lambda}$ el número de puntos críticos de índice $\lambda$ de $f$. Entonces,
    \begin{equation*}
        b_{\lambda}(M) \le C_{\lambda},~\forall \lambda\in\Lambda ~~\text{ y }~~ \chi(M) = \sum_{\lambda\in\Lambda} (-1)^{\lambda} C_{\lambda},
    \end{equation*}
    donde $\Lambda$ incluye todos los valores de los índices de los puntos críticos de $f$.
\end{theorem}
\begin{proof}
    
\end{proof}

\bigskip

\begin{theorem}[Desigualdades de Morse]

\end{theorem}
\begin{proof}
    
\end{proof}

lo de 'as an ilustration'\dots

\begin{corollary}
    Si $C_{\lambda+1} = C_{\lambda-1} = 0$ entonces $b_{\lambda} = C_{\lambda}$ y $b_{\lambda+1} = b_{\lambda-1} = 0$ 
\end{corollary}

explicar por que esto es la ostia y por que se deduce del ultimo de homotopia y de las desigualdades de morse. mirar lo del proyectivo o yo que se.